% ============================================================================
\section{Workloads, infrastructure et mesures}
\label{sec:mesures}
% ============================================================================

Avant d'optimiser, il faut savoir combien de temps prend chaque tâche sur chaque
machine. Toutes les valeurs de cette section sont des \mesure{}s.

\subsection{Workloads}

\begin{table}[H]
\centering\small
\begin{tabular}{llll}
\toprule
\textbf{Workload} & \textbf{Profil} & \textbf{Signification de \code{-{}-size}} & \textbf{Tailles mesurées} \\
\midrule
\code{cpu\_task}    & CPU          & plus grand entier testé (nombres premiers) & 200\,000 ; 500\,000 ; 1\,000\,000 \\
\code{memory\_task} & mémoire      & mémoire allouée et parcourue (Mo)          & 128 ; 256 ; 512 \\
\code{matrix\_task} & CPU + mémoire& dimension $N$ des matrices $N\times N$      & 200 ; 400 ; 600 \\
\code{sort\_task}   & mixte        & nombre d'entiers triés                    & 1\,M ; 5\,M ; 10\,M \\
\bottomrule
\end{tabular}
\caption{Les quatre workloads de référence et leurs trois tailles (celles de l'énoncé).}
\end{table}

Chaque programme affiche une ligne JSON contenant son propre temps de calcul
(\code{runtime\_s}), ce qui permet plus tard de séparer le calcul du surcoût
d'orchestration. Les tailles mesurées sont exactement celles utilisées dans les
lots de tâches : le scheduler lit directement une durée mesurée, sans extrapoler.

\subsection{Infrastructure}

\begin{table}[H]
\centering\small
\begin{tabular}{lllll}
\toprule
\textbf{VM} & \textbf{Type GCP} & \textbf{vCPU} & \textbf{RAM} & \textbf{Prix (\$/h)} \\
\midrule
cloudopt-vm1     & e2-standard-2 & 2            & 8 Go & 0,0737 \\
cloudopt-vm2     & e2-medium     & 2 (partagés) & 4 Go & 0,0369 \\
cloudopt-vm3     & e2-medium     & 2 (partagés) & 4 Go & 0,0369 \\
cloudopt-highcpu & e2-highcpu-4  & 4            & 4 Go & 0,1088 \\
\bottomrule
\end{tabular}
\caption{Les quatre VM (Ubuntu 24.04, x86\_64, zone europe-west1-b). Facturation GCP : à la seconde, minimum 60~s par démarrage.}
\label{tab:vms}
\end{table}

\subsection{Benchmark : une tâche seule}

Protocole (\code{scripts/benchmark.sh}, identique en local et sur chaque VM) :
4 workloads $\times$ 3 tailles $\times$ 5 répétitions, une exécution de chauffe
non mesurée avant chaque taille. Nous retenons la \textbf{médiane}, plus robuste
qu'une moyenne à une exécution perturbée. L'écart-type est inférieur à 1\,\%
dans presque tous les cas.

\begin{table}[H]
\centering\small
\begin{tabular}{llrrrrr}
\toprule
\textbf{Workload} & \textbf{Taille} & \textbf{Local} & \textbf{vm1} & \textbf{vm2} & \textbf{vm3} & \textbf{highcpu} \\
\midrule
cpu\_task    & 200\,000   & 0,003 & 0,045 & 0,047 & 0,045 & 0,045 \\
             & 500\,000   & 0,008 & 0,148 & 0,150 & 0,147 & 0,148 \\
             & 1\,000\,000& 0,019 & 0,378 & 0,383 & 0,378 & 0,379 \\
matrix\_task & 200        & 0,004 & 0,018 & 0,020 & 0,019 & 0,018 \\
             & 400        & 0,023 & 0,101 & 0,105 & 0,101 & 0,104 \\
             & 600        & 0,071 & 0,320 & 0,328 & 0,317 & 0,318 \\
memory\_task & 128 Mo     & 0,039 & 0,206 & 0,224 & 0,207 & 0,222 \\
             & 256 Mo     & 0,078 & 0,406 & 0,435 & 0,402 & 0,427 \\
             & 512 Mo     & 0,149 & 0,797 & 0,854 & 0,793 & 0,846 \\
sort\_task   & 1\,M       & 0,055 & 0,137 & 0,140 & 0,138 & 0,138 \\
             & 5\,M       & 0,294 & 0,720 & 0,733 & 0,719 & 0,726 \\
             & 10\,M      & 0,612 & 1,498 & 1,515 & 1,496 & 1,507 \\
\bottomrule
\end{tabular}
\caption{Médiane du temps d'une tâche seule, en secondes (\code{data/vm\_performance.csv}). Local = VM Linux OrbStack sur Mac Apple Silicon.}
\label{tab:bench}
\end{table}

\paragraph{Constats.}
\begin{enumerate}
  \item \textbf{Les quatre VM ont la même vitesse pour une tâche seule}, à 5\,\%
        près. Un workload n'utilise qu'un cœur, et toutes les VM e2 tournent sur la
        même famille de processeurs. HighCPU n'a pas des cœurs plus rapides : elle
        en a plus. Cela contredit notre première version, où HighCPU semblait
        dominer les autres VM tâche par tâche.
  \item La machine locale est 2,5 à 20 fois plus rapide : elle sert à caractériser
        les workloads, pas à prédire les durées sur GCP.
  \item Sur les VM, une tâche dure de 0,02 à 1,5~s : le surcoût d'orchestration
        n'est pas négligeable.
\end{enumerate}

\subsection{Parallélisme réel : l'expérience « workers »}
\label{sec:workers}

Le benchmark mesure une tâche \emph{seule}. Or l'orchestrateur lance plusieurs
tâches en même temps sur une VM. Nous avons donc mesuré le parallélisme réel
(\code{scripts/exp\_workers.sh}) : le même lot de 24 tâches (2 $\times$ les 12
combinaisons) sur chaque VM, avec 1, 2 puis 4 tâches simultanées
(\emph{workers}), 3 répétitions par niveau.

\begin{figure}[H]
\centering
\begin{tikzpicture}
\begin{axis}[
  width=0.62\textwidth, height=6cm,
  xlabel={Workers (tâches simultanées par VM)}, ylabel={Débit (tâches/s)},
  xtick={1,2,4}, ymin=0, ymax=6, grid=major, grid style={gray!20},
  legend style={font=\small, at={(0.5,-0.2)}, anchor=north, legend columns=2,
                /tikz/every even column/.append style={column sep=0.6cm}}]
  \addplot[mark=*, thick, selcol] coordinates {(1,2.22) (2,3.98) (4,5.32)};
  \addlegendentry{highcpu (4 vCPU)}
  \addplot[mark=square*, thick, accent] coordinates {(1,1.96) (2,2.66) (4,2.72)};
  \addlegendentry{vm1 (2 vCPU)}
  \addplot[mark=triangle*, thick, rrcol] coordinates {(1,2.19) (2,2.97) (4,3.08)};
  \addlegendentry{vm2 (2 vCPU partagés)}
  \addplot[mark=diamond*, thick, gray!50!black, dashed] coordinates {(1,2.22) (2,2.99) (4,2.86)};
  \addlegendentry{vm3 (2 vCPU partagés)}
\end{axis}
\end{tikzpicture}
\caption{Débit mesuré selon le nombre de tâches simultanées (moyenne de 3 répétitions, écart < 3\,\%).}
\label{fig:workers}
\end{figure}

\paragraph{Constats.}
\begin{enumerate}
  \item \textbf{Une VM à 2~vCPU ne vaut qu'environ 1,35 cœur} : passer de 1 à 2
        workers n'augmente le débit que de 35\,\%, car chaque tâche ralentit
        d'environ $\times 1{,}4$. Avec 4 workers, plus aucun gain. Sur GCP, un vCPU
        est un \emph{hyperthread} : 2~vCPU correspondent à un seul cœur physique.
  \item \textbf{HighCPU (4 vCPU, soit 2 cœurs) passe bien à l'échelle} :
        $\times 1{,}8$ avec 2 workers, $\times 2{,}4$ avec 4.
  \item Nous retenons \textbf{workers = nombre de vCPU} (2, ou 4 pour highcpu) :
        c'est le débit maximal ou quasi maximal partout.
\end{enumerate}

\begin{table}[H]
\centering\small
\begin{tabular}{lrrr}
\toprule
\textbf{VM} & \textbf{Débit (tâches/s)} & \textbf{Prix (\$/h)} & \textbf{Débit par \$/h} \\
\midrule
vm2 / vm3 (e2-medium)     & 2,97 / 2,99 & 0,0369 & \textbf{$\approx$ 81} \\
highcpu (e2-highcpu-4)    & 5,32        & 0,1088 & 49 \\
vm1 (e2-standard-2)       & 2,66        & 0,0737 & 36 \\
\bottomrule
\end{tabular}
\caption{Rapport débit/prix au nombre de workers retenu. Les e2-medium sont les plus économiques, highcpu la plus rapide, vm1 est dominée.}
\label{tab:debitprix}
\end{table}

Ce tableau résume le compromis que doit arbitrer le scheduler : \textbf{vm2/vm3
pour économiser, highcpu pour aller vite, vm1 seulement si la capacité manque}.
